\subsection*{11.3 Application: Monte Carlo Integration}

Suppose we have a multi-variable function g(x1,x 2,...,x d) with d variables,

where each variable takes a value between 0 and 1. Our goal is to evaluate the

integral

\int 1

\int 1

\cdot\cdot\cdot

\int 1

g(x1,x 2,...,x d)dx 1dx2 \cdot\cdot\cdot dxd, (11.2)

which is a d -fold integration over the unit hypercube. A direct method that

approximates this integral by a Riemann sum involves selecting m points along

each dimension and computing the sum

md

m- 1\sum

j1=0

m- 1\sum

j2=0

\cdot\cdot\cdot

m- 1\sum

jd=0

g

( j1

m, j2

m,..., jd

m

)

However, this method requires evaluating the function at md points, which can

be computationally infeasible when d is large. This is known as the "curse of

dimensionality", where even a small number of sample points per dimension can

result in an exponential increase in the total number of function evaluations required.

To overcome this limitation, we can use Monte Carlo integration, which involves

randomly sampling points from the domain of the hypercube .[0,1]d , and taking

the average of the function values at the random points. To estimate the integral

using n sample points, we generate independent uniform random variables Xi1 ,

Xi2,...,X id between 0 and 1, for i= 1 ,2,...,n We then compute the estimator

In \coloneqq 1

n

n\sum

i=1

g(Xi1,Xi2,...,X id ). (11.3)

The n terms in this estimator are independent, and its expected value is given by

E[In]= 1

n

n\sum

i=1

E[g(Xi1,Xi2,...,X id )]= E[g(X11,X 12,...,X 1d)],

which is equal to the multi-variable integral in (11.2)If the integral in (11.2) is

finite, then by Khinchin's weak law of large numbers, the estimator In converges in

probability to its expected value E[In] asn\to\infty .

The convergence rate of the Monte Carlo estimator is independent of the

dimensionality of the problem. To see why, suppose that the function g to be

integrated is bounded by a constant C in absolute value, and suppose the error

tolerance is. \epsilon. Let the true value of the integral be I*Because the variance is finite,

by the. L2 version of the weak law of large numbers, we have

P

( \vert\vert\vertIn -I *

\vert\vert\vert >\epsilon

)

\leq Va r(g(X 11,...,X 1d))

n\epsilon2 \leq C2

n\epsilon2 .

If we want the probability P( \vertIn -I *\vert>\epsilon ) to be less than, say, 0.01, we need

to take at least C2/(0.01\epsilon2) samples. The bound increases inversely proportional

to. \epsilon2. This error estimate does not depend on the number of variables and is valid in

any dimension. This means that the convergence rate of the Monte Carlo estimator

is the same for one-dimensional problems as it is for high-dimensional problems.

However, the variance of the estimator can still be high for functions with large

fluctuations, requiring the evaluation of the function on many sample points to

achieve a certain level of accuracy.

We illustrate the convergence speed of Monte Carlo integration with the follow-

ing example.

Python Program for Monte Carlo Integration We demonstrate the Monte Carlo

method by evaluating the integral

\int M

-M

\int M

-M

\int M

-M

\int M

-M

x2y2z2w2 dxdydzdw

with four variables. The integral is separable, with exact value (2M3/3)4W e

can implement the Monte Carlo integration method using the following Python

program:

from random import uniform

n = 100000 # number of sample points

f = lambda u: (u[0] *u[1]*u[2]*u[3])**2 # function to be integrated

M = 2 # range of the integration

s=0

for i in range(n):

random_point = [uniform(-M,M) for j in range(4)]

s += f(random_point) # evaluate function f

print("Integral = ", (2 *M)**4*s/n)

We set M = 2. The exact value of the integral is 809.086. The Python function

uniform(-M,M) generates a float-type number uniformly between -M and MW e

run the program with n = 100, 103, 104, 105, and 10 6. The approximate values

produced by the program are shown in the table below.

n 100 1000 10000 100000 1000000

In 724.062 822.986 796.259 812.225 807.883

When n is large, it is close to the exact value 809.086.

Quasi-Random Monte Carlo

Another method for numerical evaluation of multi-variable is to use quasi-random point set instead

of pseudo-random numbers. Quasi-random point sets are sequence of points that are evenly spread

in the hypercube .[0, 1]d but are not statistically the same as uniformly random points. These

sequences are also called "low-discrepancy sequences"If we use quasi-random point sets in Monte

Carlo integration, the estimate can converge to the true value of the integral faster than random

points, especially in high-dimensional problems.
